\documentclass[10pt]{article}
\usepackage[T1]{fontenc}
\usepackage[utf8]{inputenc}
\usepackage[french]{babel}
\usepackage{amsmath,amssymb,mathrsfs,mathpazo}
\usepackage{geometry}
\geometry{a4paper,margin=2cm}
\usepackage{graphicx,array,tabularx,tabularray,colortbl,multirow,multicol,booktabs,enumitem,pifont,ulem,fancybox,diagbox,cancel}
\usepackage{pstricks,pstricks-add,pst-plot,pst-tree,pst-node,pst-text,pst-func,pst-math,pst-bspline,pst-3dplot,pst-eucl,pst-solides3d}
\let\pscalculate\undefined
\usepackage{pst-bezier}
\usepackage{tikz,pgfplots,tkz-tab}
\usetikzlibrary{arrows,arrows.meta,patterns,calc,positioning,trees,decorations.markings,shapes.geometric,shapes.misc,shapes.arrows,intersections,decorations.pathmorphing,decorations.pathreplacing,backgrounds,shadows,babel}
\pgfplotsset{compat=1.18}
\usepgfplotslibrary{fillbetween}
\usepackage[np]{numprint}
\usepackage{listings,makecell,pas-tableur,twoopt}
\let\barre\undefined
\usepackage[pstricks]{tablvar}
\definecolor{darkgreen}{RGB}{0,100,0}
\providecommand{\red}{\color{red}}
\providecommand{\blue}{\color{blue}}
\providecommand{\green}{\color{green}}
\providecommand{\black}{\color{black}}
\providecommand{\white}{\color{white}}
\providecommand{\gray}{\color{gray}}
\providecommand{\cyan}{\color{cyan}}
\providecommand{\magenta}{\color{magenta}}
\providecommand{\yellow}{\color{yellow}}
\providecommand{\violet}{\color{violet}}
\providecommand{\orange}{\color{orange}}
\providecommand{\euro}{\texteuro}
\providecommand{\EUR}{\texteuro}
\providecommand{\ieme}{\textsuperscript{e}}
\providecommand{\degres}{\ensuremath{^\circ}}
\providecommand{\degree}{\ensuremath{^\circ}}
\providecommand{\vv}[1]{\overrightarrow{#1}}
\providecommand{\Coord}[1]{#1}
\providecommand{\decosix}{}
\providecommand{\starredbullet}{\ensuremath{\bullet}}
\pagestyle{empty}
\setlength{\parindent}{0pt}
\usepackage{adjustbox}\DeclareUnicodeCharacter{2236}{:}\DeclareUnicodeCharacter{2190}{\ensuremath{\leftarrow}}
\providecommand{\R}{\mathbb{R}}
\providecommand{\N}{\mathbb{N}}
\providecommand{\Z}{\mathbb{Z}}
\providecommand{\Q}{\mathbb{Q}}
\providecommand{\D}{\mathbb{D}}
\providecommand{\C}{\mathbb{C}}
\providecommand{\e}{{\mathrm{e}}}
\providecommand{\vect}[1]{\overrightarrow{#1}}
\providecommand{\vectt}[1]{\overrightarrow{\mathrm{#1}}}
\providecommand{\vvt}[1]{\overrightarrow{\mathrm{#1}}}
\providecommand{\np}[2][]{{#2~#1}}
\providecommand{\vv}[1]{{\overrightarrow{#1}}}
\renewcommand{\d}{{\mathrm{d}}}
\providecommand{\degres}{{^\circ}}
\providecommand{\degre}{{^\circ}}
\providecommand{\blue}{\color{blue}}
\providecommand{\red}{\color{red}}
\providecommand{\enskip}{\quad}
\providecommand{\llbracket}{\lbrack\!\lbrack}
\providecommand{\rrbracket}{\rbrack\!\rbrack}
\providecommand{\euro}{{\text{€}}}
\providecommand{\no}{{\text{n°}}}
\providecommand{\interval}[2]{{[#1;#2]}}
\providecommand{\intervFO}[2]{{[#1;#2[}}
\providecommand{\dotfill}{\ldots}

\providecommand{\ds}{\displaystyle}
\providecommand{\textding}[1]{#1}
\providecommand{\Oij}{\(\mathrm{O};\vec{\imath},\vec{\jmath}\)}
\providecommand{\Oijk}{\(\mathrm{O};\vec{\imath},\vec{\jmath},\vec{k}\)}
\providecommand{\Ouv}{\(\mathrm{O};\vec{u},\vec{v}\)}
\providecommand{\intervFF}[2]{\ensuremath{[#1;#2]}}
\renewcommand{\euro}{\ifmmode\text{\texteuro}\else\texteuro\fi}
\DeclareUnicodeCharacter{00B0}{\ifmmode^{\circ}\else\textdegree\fi}
\DeclareUnicodeCharacter{20AC}{\ifmmode\text{\texteuro}\else\texteuro\fi}
% Énoncés : Éducation nationale. Transcription et corrigés : APMEP, Denis Vergès et contributeurs. Découpage et classement : CoopMaths, Sébastien Lozano, Éric Elter et contributeurs. Adaptation de la présentation : Maths Explorer.
% CC BY-SA : https://coopmaths.fr/www/about/
% https://coopmaths.fr/alea/static/eam/2026/tex/eam_specialite_2026_09_metropole_2_cor.tex
\begin{document}
\noindent\textit{APMEP / CoopMaths. CC BY-SA. Présentation adaptée par Maths Explorer.}\par\medskip
\emph{Pour chaque affirmation, on indique si elle est vraie ou fausse en justifiant ; une
réponse non justifiée n'est pas prise en compte. Les trois questions sont indépendantes.}

\begin{enumerate}
	\item \textbf{Affirmation fausse.}

	Pour tout entier naturel $n$, avec $v_n = 3 \times \left(\dfrac{1}{4}\right)^n$ :
	\[v_{n+1} - v_n = 3 \times \left(\dfrac{1}{4}\right)^{n+1} - 3 \times \left(\dfrac{1}{4}\right)^n
	= 3 \times \left(\dfrac{1}{4}\right)^n \times \left(\dfrac{1}{4} - 1\right)
	= -\dfrac{9}{4} \times \left(\dfrac{1}{4}\right)^n.\]
	Or $\left(\dfrac{1}{4}\right)^n > 0$ et $-\dfrac{9}{4} < 0$, donc $v_{n+1} - v_n < 0$ : la suite est strictement
	décroissante. 

	$(v_n)$ est une suite géométrique de premier terme $3 > 0$ et de
	raison $\dfrac{1}{4} \in \,]0 ; 1[$ : elle est \textbf{décroissante}.

	\item \textbf{Affirmation vraie.}

	On se ramène à une équation du second degré :
	\[x(1-x) = 4 \iff x - x^2 = 4 \iff -x^2 + x - 4 = 0 \iff x^2 - x + 4 = 0.\]
	$\Delta = (-1)^2 - 4 \times 1 \times 4 = 1 - 16 = -15 < 0$.

	Comme $\Delta < 0$, l'équation $x^2 - x + 4 = 0$ n'a aucune
	solution réelle. 
	
	Aucun réel $x$ ne vérifie $x(1-x) = 4$.

	\item \textbf{Affirmation vraie.}

	
D'après la relation de Chasles :
\[
\overrightarrow{KD}=\overrightarrow{KA}+\overrightarrow{AD}
\qquad\text{et}\qquad
\overrightarrow{KL}=\overrightarrow{KB}+\overrightarrow{BL}.
\]
Par bilinéarité du produit scalaire :
\begin{align*}
\overrightarrow{KD}\cdot\overrightarrow{KL}
&=\left(\overrightarrow{KA}+\overrightarrow{AD}\right)\cdot\left(\overrightarrow{KB}+\overrightarrow{BL}\right)\\
&=\overrightarrow{KA}\cdot\overrightarrow{KB}
+\overrightarrow{KA}\cdot\overrightarrow{BL}
+\overrightarrow{AD}\cdot\overrightarrow{KB}
+\overrightarrow{AD}\cdot\overrightarrow{BL}.
\end{align*}


\begin{itemize}
  \item $\overrightarrow{KA}$ et $\overrightarrow{KB}$ sont colinéaires et de sens contraires (car $K\in[AB]$) :
        $\overrightarrow{KA}\cdot\overrightarrow{KB}=-KA\times KB=-3\times 1=-3$.
  \item $(AB)\perp(BC)$ et $(AD)\perp(AB)$ car $ABCD$ est un rectangle, donc
        $\overrightarrow{KA}\cdot\overrightarrow{BL}=0$ et $\overrightarrow{AD}\cdot\overrightarrow{KB}=0$.
  \item $\overrightarrow{AD}$ et $\overrightarrow{BL}$ sont colinéaires et de même sens :
        $\overrightarrow{AD}\cdot\overrightarrow{BL}=AD\times BL=3\times 2=6$.
\end{itemize}


\[
\overrightarrow{KD}\cdot\overrightarrow{KL}=-3+0+0+6=3.
\]
\end{enumerate}

\end{document}
